% =====================================================================
%  tipos_algoritmos.tex
%  ====================
%
%
%  Description:
%  ------------
%  Reporte de investigación correspondiente a la Actividad 3 del curso
%  propedéutico de Programación y Estructuras de Datos (MC en Ciencias
%  de la Computación, INAOE). Recorre las principales familias de
%  algoritmos: determinísticos, no determinísticos, divide y vencerás,
%  voraces, programación dinámica, backtracking, ramificación y
%  acotación, probabilísticos, paralelos y metaheurísticas.
%
%
%  Metadata:
%  ----------
%   * Author : zxxz6 (Bryan Violante Arriaga)
%   * Version: 1.0.0
%   * License: 
%
%
%  History:
%  ------------
%  Author      Date            Description
%  zxxz6       16/05/2026      Creación
%
%
% =====================================================================
\documentclass[12pt,letterpaper]{article}

% ---------- Codificación e idioma ----------
\usepackage[utf8]{inputenc}
\usepackage[T1]{fontenc}
\usepackage[spanish,es-tabla,es-noindentfirst]{babel}
\usepackage{lmodern}

% ---------- Geometría y diseño ----------
\usepackage[letterpaper,margin=2.5cm]{geometry}
\usepackage{setspace}
\onehalfspacing
\usepackage{microtype}
\usepackage{parskip}

% ---------- Formato del abstract  ----------
\usepackage{abstract}
\renewcommand{\abstractnamefont}{\normalfont\bfseries}
\renewcommand{\abstracttextfont}{\normalfont\small\itshape}
\setlength{\absleftindent}{0.5cm}
\setlength{\absrightindent}{0.5cm}

% ---------- Matemáticas ----------
\usepackage{amsmath,amssymb,mathtools}

% ---------- Tablas, color y captions ----------
\usepackage{graphicx}
\usepackage{booktabs}
\usepackage{array}
\usepackage{tabularx}
\usepackage{xcolor}
\usepackage{caption}
\captionsetup{font=small,labelfont=bf}
\usepackage{pdflscape}
\usepackage[section]{placeins}

% ---------- Pseudocódigo ----------
\usepackage{algorithm}
\usepackage{algpseudocode}
\floatname{algorithm}{Algoritmo}
\algrenewcommand\algorithmicwhile{\textbf{mientras}}
\algrenewcommand\algorithmicdo{\textbf{hacer}}
\algrenewcommand\algorithmicend{\textbf{fin}}
\algrenewcommand\algorithmicif{\textbf{si}}
\algrenewcommand\algorithmicthen{\textbf{entonces}}
\algrenewcommand\algorithmicelse{\textbf{si no}}
\algrenewcommand\algorithmicfor{\textbf{para}}
\algrenewcommand\algorithmicforall{\textbf{para todo}}
\algrenewcommand\algorithmicreturn{\textbf{devolver}}
\algrenewcommand\algorithmicrequire{\textbf{Entrada:}}
\algrenewcommand\algorithmicensure{\textbf{Salida:}}
\algrenewcommand\algorithmicrepeat{\textbf{repetir}}
\algrenewcommand\algorithmicuntil{\textbf{hasta que}}

% ---------- Hipervínculos ----------
\usepackage{hyperref}
\hypersetup{
    colorlinks=true,
    linkcolor=black,
    urlcolor=blue!50!black,
    pdftitle={Un viaje breve por los tipos de algoritmos},
    pdfauthor={Bryan Ezequiel Violante Arriaga}
}

% ---------- Portada ----------
\title{\textbf{Un viaje breve por los tipos de algoritmos}}
\author{Bryan Ezequiel Violante Arriaga\\
\small Prope. Programación y Estructuras de Datos\\}
\date{\today}

% =====================================================================
\begin{document}
% =====================================================================

\maketitle
\thispagestyle{empty}

% --------------------------- RESUMEN ---------------------------------
\begin{abstract}
    \noindent
    El presente trabajo (correspondiente a la actividad 3 del curso
    propedéutico de Programación y Estructuras de Datos para la MC. en Ciencias
    de la Computación) recorre los distintos tipos de algoritmos que suelen
    estudiarse en un curso de análisis y diseño. Tras unas notas sobre la
    noción de algoritmo y la medición de su eficiencia mediante notación
    asintótica, se abordan las familias clásicas: determinísticos y no
    determinísticos, divide y vencerás, voraces, programación dinámica,
    vuelta atrás (\emph{backtracking}), ramificación y acotación,
    probabilísticos, paralelos y metaheurísticas. Para cada una se
    comenta su idea principal, los problemas en los que conviene aplicarla
    y, cuando resulta útil, un pequeño ejemplo en pseudocódigo. Al cierre
    se incluye una tabla comparativa y algunas reflexiones sobre cuándo
    elegir una técnica u otra.
    \vspace{0.5em}
\end{abstract}


% --------------------------- INTRODUCCIÓN ----------------------------
\section{Introducción}
\label{sec:introduccion}

Resolver problemas con una computadora implica, básicamente, diseñar un 
algoritmo, que en palabras simples es una receta paso a paso que recibe una 
entrada y produce una salida (similar a una receta de cocina, por ejemplo). 
Tareas como ordenar una lista, calcular la ruta más corta entre dos ciudades, 
decidir si un número es primo o planificar tareas en una fábrica son ejemplos 
cotidianos para un científico de la computación. La buena noticia es que, con 
el paso del tiempo, las personas dedicadas a esto han identificado varios 
\emph{patrones} de diseño que resuelven familias enteras de problemas. 
La mala noticia es que no hay un único patrón que sirva para todo.

En este trabajo se recorren los patrones más conocidos y trata de responder, de 
manera práctica y concisa, tres preguntas para cada uno: ¿en qué consiste?, 
¿qué tipo de problemas resuelve bien? y ¿cuánto cuesta usarlo?

% --------------------------- MARCO TEÓRICO ---------------------------
\section{Marco teórico}
\label{sec:marco}

\subsection{¿Qué es un algoritmo?}
Como se mencionó, se puede pensar en un algoritmo como una receta de cocina: más
formalmente, un algoritmo es una secuencia finita de instrucciones bien
definidas que transforma una entrada en una salida. Las propiedades clave que
suele exigirse son cinco:

\begin{itemize}
    \item \textbf{Finitud:} termina en un número finito de pasos.
    \item \textbf{Definitud:} cada paso es claro y sin ambigüedad.
    \item \textbf{Entrada:} recibe cero o más datos.
    \item \textbf{Salida:} produce uno o más resultados.
    \item \textbf{Efectividad:} cada operación es realizable.
\end{itemize}

\subsection{Complejidad}
La eficiencia de un algoritmo se mide observando cuántas operaciones realiza en 
función del tamaño $n$ de la entrada. A esta cuenta se le llama 
\emph{complejidad temporal}, $T(n)$. De manera similar, la 
\emph{complejidad espacial} mide la memoria que necesita.

\subsection{Notación asintótica}
Para describir el crecimiento de $T(n)$ se usan tres notaciones:

\begin{itemize}
    \item \textbf{Cota superior:} $f(n) = \mathcal{O}(g(n))$ si existen 
    constantes $c>0$ y $n_0$ tales que
    \begin{equation}
        0 \leq f(n) \leq c\,g(n), \quad \forall n \geq n_0.
    \end{equation}
    \item \textbf{Cota inferior:} $f(n) = \Omega(g(n))$ si existen $c>0$ y 
    $n_0$ tales que
    \begin{equation}
        0 \leq c\,g(n) \leq f(n), \quad \forall n \geq n_0.
    \end{equation}
    \item \textbf{Cota ajustada:} $f(n) = \Theta(g(n))$ si se cumplen ambas a 
    la vez.
\end{itemize}

Un algoritmo es \emph{polinómico} si $T(n) = \mathcal{O}(n^k)$ para alguna 
constante $k$, y \emph{exponencial} si $T(n) = \mathcal{O}(c^n)$ con $c>1$. 
A simple vista (al menos para el autor) resultaría  que no hay mucha diferencia, 
pero esto no es así, ya que la diferencia es enorme en la práctica.\\
A continuación se muestra una tabla comparativa, en ella se destaca el tiempo de
ejecución para diferentes tamaños de entrada y diferentes tiempos de ejecución, 
suponiendo que el procesador puede realizar $10^9$ operaciones por segundo 
(siendo modestos pues es común que los equipos de hoy en día superen esa cifra).

\begin{table}[ht]
\centering
\small
\caption{Tiempo de ejecución (orden de magnitud) suponiendo
$10^{9}$ operaciones por segundo.}
\label{tab:explosion}
\begin{tabular}{|l|c|c|c|}
\hline
$T(n)$ & $n=10$ & $n=30$ & $n=50$ \\
\hline
$n$    & $\sim$10\,ns      & $\sim$30\,ns   & $\sim$50\,ns        \\
$n^2$  & $\sim$0.1\,$\mu$s & $\sim$1\,$\mu$s & $\sim$2.5\,$\mu$s   \\
$n^5$  & 0.1\,ms           & $\sim$24.3\,ms    & $\sim$0.3\,s      \\
$2^n$  & 1\,$\mu$s         & $\sim$1\,s     & $\approx$13\,días   \\
$3^n$  & 59\,$\mu$s        & $\approx$57\,hrs    & $\approx$22\,millones de años    \\
\hline
\end{tabular}
\end{table}

\subsection{Clases de problemas}
De manera informal, un problema está en la clase \textsf{P} si existe un 
algoritmo polinómico que lo resuelve, y está en \textsf{NP} si una solución 
propuesta puede verificarse en tiempo polinómico. Los problemas 
\textsf{NP}-duros son los más complejos dentro de \textsf{NP}. Si estos 
problemas tienen o no algoritmos polinómicos es la famosa pregunta abierta 
$\mathsf{P}\stackrel{?}{=}\mathsf{NP}$.

% --------------------------- TIPOS DE ALGO ------------------------------


% ---------- Determinísticos ----------
\section{Algoritmos determinísticos}
\label{sec:det}

Un algoritmo es \emph{determinístico} cuando, para una misma entrada, su 
comportamiento es siempre el mismo: cada paso lleva a un único siguiente paso 
y la salida está completamente determinada por la entrada. Es el tipo de 
algoritmo más común y, en cierto sentido, el ``predeterminado'' que pensamos al
momento de pensar en un algoritmo para resolver un problema específico; hacer A, 
luego B, y así sucesivamente.

Este tipo de algoritmo puede servir como solución o enfoque correcto para 
muchos problemas, prácticamente todos los algoritmos clásicos que se estudian 
al inicio entran aquí (búsqueda secuencial, ordenamiento por inserción, 
búsqueda en profundidad, etc.). Aparece sobre todo por contraste con las otras 
categorías ya que en la mayoría de los casos, la solución otorgada por este tipo 
de algoritmos no es la mejor posible.

La ejecución de este tipo de algoritmos corresponde a una función 
$f : \mathcal{I} \rightarrow \mathcal{O}$ totalmente definida. En palabras 
simples, si se ejecuta dos veces con la misma entrada, se obtiene exactamente 
la misma salida y se pasa por los mismos pasos (de manera similar a un AFD).

Un ejemplo clásico de algoritmo determinístico es la búsqueda secuencial, 
que recorre un arreglo de izquierda a derecha hasta encontrar el valor buscado. 
En el peor caso, recorre todo el arreglo, lo que da una complejidad de 
$\Theta(n)$. El pseudocódigo es el siguiente:

\begin{verbatim}
    busqueda_secuencial(A, x):
        # A es un arreglo de tamaño n, x el valor buscado
        para i = 1, 2, ..., n:
            si A[i] == x:
                devolver i
        devolver -1   # no se encontro
\end{verbatim}

Puede observarse que el algoritmo busca de manera ingenua, sin aprovechar 
ninguna propiedad de la entrada, esto lo hace poco eficiente para arreglos 
grandes, y por supuesto existen mejores algoritmos para este tipo de problemas,
los cuales se estudiarán más adelante.

% ---------- No determinísticos ----------
\section{Algoritmos no determinísticos}
\label{sec:nodet}

Un algoritmo \emph{no determinístico} se ejecuta de manera similar a un árbol: 
en cada paso el cómputo puede bifurcarse en varios sucesores y todas las ramas 
se ejecutan, en teoría, al mismo tiempo. 
Se considera que el algoritmo acepta si \emph{alguna} de las ramas llega a un 
estado de aceptación (resulta una analogía muy interesante con los AFN y la 
conversión a AFD pues en este proceso se exploran todas las posibles 
transiciones, formando un árbol).

Este tipo de algoritmos permiten definir formalmente la clase \textsf{NP}: un 
problema está en \textsf{NP} si existe una máquina no determinística que lo 
resuelve en tiempo polinómico. Cuando hay que ``simular'' este comportamiento 
en una computadora normal, lo que se hace es recorrer ese árbol de 
posibilidades con técnicas como \emph{backtracking} o ramificación y acotación,
las cuales se estudiarán en las secciones posteriores.

El principio de este algoritmo parecería simple, pero es muy poderoso, en cada 
paso, en lugar de elegir una transición, el algoritmo ``adivina'' la correcta. 
La pregunta abierta más importante de la teoría de la computación, 
$\mathsf{P}=\mathsf{NP}$, pregunta si toda esa magia adivinatoria puede 
sustituirse por un algoritmo determinístico polinómico.

% ---------- Divide y Vencerás ----------
\section{Divide y vencerás}
\label{sec:dyv}

En los anales de la historia, se le atribuye al general romano Julio César la 
frase ``divide y vencerás'', la cual consistía en dividir las fuerzas enemigas 
en concentraciones más pequeñas para derrotarlas de manera más fácil.

Esta técnica 
se basa en ese concepto antiquísimo y consiste en, dado un problema, partirlo 
en pedazos más pequeños del mismo tipo, resolver cada uno por separado 
(normalmente con la misma técnica, de forma recursiva) y combinar los 
resultados. Los tres pasos clásicos son: \textbf{dividir}, \textbf{vencer} 
(resolver los subproblemas) y \textbf{combinar}.

Funciona muy bien cuando el problema admite una partición natural en partes 
independientes y de tamaño parecido. Los ejemplos clásicos incluyen merge sort
(\emph{merge sort}), búsqueda binaria, multiplicación 
rápida de enteros grandes y el algoritmo de Strassen para multiplicar matrices.

Si un algoritmo divide y vencerás genera $a$ subproblemas de tamaño $n/b$ y 
emplea $f(n)$ tiempo en dividir y combinar, su tiempo de ejecución satisface
\begin{equation}
    T(n) = a\,T\!\left(\frac{n}{b}\right) + f(n).
\end{equation}
El \emph{teorema maestro} permite resolver estas recurrencias de manera
sistemática (de manera muy simplificada, este teorema dice que basta con 
comparar $f(n)$ con $n^{\log_b a}$, el costo acumulado de los subproblemas y, 
según cuál de los dos domine, entrega directamente la forma cerrada de $T(n)$ 
sin necesidad de desarrollar el árbol de recursión a mano).

A continuación se muestra el pseudocódigo del algoritmo de merge sort, 
que divide el arreglo a ordenar en dos mitades, las ordena recursivamente y 
luego las fusiona en un arreglo ordenado. La función \textsc{merge} se encarga 
de combinar las dos mitades ya ordenadas, y su complejidad es $\Theta(n)$.

\begin{verbatim}
    merge_sort(A, p, r):
        # A es el arreglo, p y r los indices del subarreglo a ordenar
        si p < r:
            q = (p + r) // 2
            merge_sort(A, p, q)
            merge_sort(A, q+1, r)
            merge(A, p, q, r)   # fusiona las dos mitades ya ordenadas
\end{verbatim}

Se puede ver que la recurrencia $T(n) = 2T(n/2) + \Theta(n)$ da $T(n) = 
\Theta(n \log n)$, que es lo mejor que puede hacerse para ordenamiento basado 
en comparaciones.

% ---------- Voraces ----------
\section{Algoritmos voraces (\emph{greedy})}
\label{sec:greedy}

En la vida diaria, al estar hambriento y, presentarse diversos platillos para 
comer, se tiende a elegir lo más apetecible, al terminar, se elige el siguiente 
platillo más apetecible de entre las opciones restantes y así sucesivamente. 
De manera similar, un algoritmo voraz construye la solución paso a paso: en 
cada momento elige el candidato que se ve más prometedor según un criterio 
local, y no se vuelve atrás. Es la estrategia del que va decidiendo 
``lo mejor que se ve ahora'', con la esperanza de llegar a un buen resultado 
global.

Este enfoque se aplica a problemas de optimización. Funciona perfectamente 
cuando el problema cumple dos propiedades: 
\textbf{subestructura óptima} (la solución óptima del problema contiene 
soluciones óptimas de sus subproblemas) y \textbf{elección voraz} 
(existe una decisión local que pertenece a alguna solución óptima).

Ejemplos clásicos donde sí funciona: árbol de recubrimiento mínimo 
(Kruskal, Prim), códigos de Huffman, camino más corto desde un origen 
(Dijkstra), selección de actividades.

Es importante aclarar que en muchos problemas la estrategia voraz NO da el 
resultado óptimo. En el problema del agente viajero, por ejemplo, ir siempre 
al vecino más cercano da rutas que pueden ser mucho peores que la mejor posible.

Un algoritmo voraz suele tener cinco partes:
\begin{itemize}
    \item El conjunto de candidatos.
    \item Los conjuntos de seleccionados y rechazados.
    \item Una función de factibilidad (¿se puede seguir armando una solución?).
    \item Una función de selección (¿cuál es el mejor candidato disponible?).
    \item Una función objetivo (¿cuándo se ha alcanzado una solución completa?).
\end{itemize}

A continuación se muestra un ejemplo de algoritmo voraz: el algoritmo de Kruskal
, que construye un árbol de recubrimiento mínimo (MST) eligiendo en cada paso 
la arista de menor peso que no forme ciclo.

\begin{verbatim}
    kruskal(G, w):
        # G = (V, E) grafo conexo, w: E -> R+ peso de cada arista
        T = {}                       # conjunto de aristas del MST
        para cada v en V:
            make_set(v)              # cada vertice en su propio componente
        ordenar E por peso w ascendente
        para cada arista (u, v) en E (en orden):
            si find_set(u) != find_set(v):
                T = T union {(u, v)} # no forma ciclo, se acepta
                union(u, v)
        devolver T
\end{verbatim}


% ---------- Programación Dinámica ----------
\section{Programación dinámica}
\label{sec:dp}

Al igual que divide y vencerás, la programación dinámica también descompone el 
problema en subproblemas, pero a diferencia de este, los subproblemas 
\emph{se repiten}. La idea entonces es resolver cada subproblema una sola vez 
y guardar el resultado en una tabla para reutilizarlo. Se cambia memoria por 
tiempo, y a veces el ahorro es enorme: un algoritmo recursivo ingenuo 
exponencial puede volverse polinómico con esta técnica.

Este tipo de algoritmos funciona muy bien cuando el problema cumple dos 
características muy específicas:
\begin{enumerate}
    \item \textbf{Subestructura óptima:} la solución óptima del problema se 
    construye con soluciones óptimas de subproblemas.
    \item \textbf{Subproblemas solapados:} los mismos subproblemas aparecen 
    una y otra vez.
\end{enumerate}
Ejemplos típicos: mochila 0-1, distancia de edición, parentización óptima de 
productos de matrices, problema del cambio de moneda, alineamiento de 
secuencias en bioinformática.

Los algoritmos basados en programación dinámica suelen seguir este 
patrón/esquema general:
\begin{enumerate}
    \item Se caracteriza cómo se ve una solución óptima.
    \item Se define recursivamente el valor de una solución óptima.
    \item Se calcula ese valor, normalmente \emph{de abajo hacia arriba} 
    (\emph{bottom-up}) aunque también se puede hacer \emph{de arriba hacia 
    abajo} (\emph{top-down}) con memoización.
    \item Se reconstruye la solución óptima a partir de la información guardada.
\end{enumerate}
Y hay dos formas equivalentes de implementarlo: la \emph{memoización} 
(recursión que guarda resultados) y la \emph{tabulación} (iterativo, llenando 
la tabla).

En la experiencia del autor, la programación dinámica es de las técnicas más 
difíciles de entender al inicio y de dominar, sobre todo por la necesidad de 
identificar la estructura de los subproblemas y la relación entre ellos. Sin 
embargo, una vez que se entiende, es una herramienta extremadamente poderosa.

A continuación se muestra un ejemplo de algoritmo de programación dinámica 
(meramente ilustrativo); el problema del cambio de moneda, el cual consiste en, 
dada una cantidad $N$ y un conjunto de denominaciones $\{d_1,\dots,d_n\}$, 
encontrar el número mínimo de monedas para pagar $N$. Sea $c[i,j]$ el número 
mínimo de monedas para pagar $j$ usando denominaciones $\{d_1,\dots,d_i\}$. 
La recurrencia es:
\begin{equation}
    c[i,j] = \min\bigl(c[i-1,j],\; 1 + c[i, j - d_i]\bigr).
\end{equation}

Y el pseudocódigo es el siguiente:

\begin{verbatim}
    cambio_monedas(d, N):
        # d = denominaciones [d_1, ..., d_n], N = cantidad objetivo
        # c[i][j] = monedas minimas para formar j usando solo d[1..i]
        para i = 1, ..., n:
            c[i][0] = 0              # 0 monedas para cantidad 0
        para i = 1, ..., n:
            para j = 1, ..., N:
                si i == 1 y j < d[i]:
                    c[i][j] = +infinito       # no alcanzable
                si no, si i == 1:
                    c[i][j] = 1 + c[i][j - d[i]]
                si no, si j < d[i]:
                    c[i][j] = c[i-1][j]       # no cabe la moneda d[i]
                si no:
                    c[i][j] = min( c[i-1][j],            # no uso d[i]
                                   1 + c[i][j - d[i]] )  # uso d[i]
        devolver c[n][N]
\end{verbatim}

La complejidad de este algoritmo es $\Theta(nN)$ (pseudopolinómica). 

% ---------- Backtracking ----------
\section{Vuelta atrás (\emph{backtracking})}
\label{sec:back}

El \emph{backtracking} explora el espacio de soluciones como si fuera un árbol: 
va construyendo una solución parcial paso a paso, y cuando descubre que esa 
rama no puede llevar a una solución válida, retrocede y prueba otra. Es 
básicamente una búsqueda en profundidad con poda por factibilidad.

Cualquier problema donde la solución se pueda expresar como una secuencia de 
decisiones $(x_1,x_2,\dots,x_k)$ con restricciones que se verifican sobre la 
marcha. Algunos ejemplos clásicos son: las $n$-reinas, suma de subconjuntos, 
coloreo de grafos, generación de permutaciones, sudoku.

A continuación se muestra el esquema general de un algoritmo de backtracking, 
que recibe una solución parcial $\mathbf{x}=(x_1,\dots,x_k)$ y el nivel $k$ del 
árbol. Si $\mathbf{x}$ es una solución completa, se procesa (por ejemplo, se 
imprime). Si no, se generan los candidatos para $x_{k+1}$ y se verifica su 
factibilidad.

\begin{verbatim}
    backtrack(x, k):
        # x = (x_1, ..., x_k) es la solucion parcial construida hasta ahora
        si x es una solucion completa:
            procesar(x)
        si no:
            para cada candidato v para x_{k+1}:
                si (x, v) es factible:
                    x_{k+1} = v
                    backtrack(x, k+1)
                    # al volver, se descarta v y se prueba otro
\end{verbatim}

En el peor caso es exponencial, pero la poda por factibilidad suele recortar 
de manera significativa el árbol. Para las 8-reinas, por ejemplo, se examinan 
cerca de 1.3 millones de configuraciones en lugar de los más de 16 millones del 
enfoque más obvio.

% ---------- Ramificación y poda ----------
\section{Ramificación y acotación (\emph{branch and bound})}
\label{sec:bnb}

La ramificación y acotación es como un \emph{backtracking} con cerebro: además 
de descartar ramas por factibilidad, calcula una \emph{cota} (superior o 
inferior, según se busque máximo o mínimo) del mejor valor posible alcanzable 
desde cada nodo. Si esa cota ya es peor que la mejor solución conocida, no 
tiene sentido seguir explorando esa rama y se poda.

Este tipo de algoritmos son muy buenos para resolver problemas de optimización 
combinatoria \textsf{NP}-duros para los que aún así se quiere la solución 
óptima exacta: 

\begin{itemize}
    \item \textbf{Mochila 0-1:} dado un conjunto de objetos con peso y valor, 
    y una capacidad máxima, se busca el subconjunto de objetos que maximice el 
    valor sin exceder la capacidad.
    \item \textbf{Programación entera:} optimización lineal con la restricción 
    de que algunas o todas las variables deben ser enteras.
    \item \textbf{Asignación cuadrática:} asignar $n$ tareas a $n$ agentes 
    minimizando un costo que depende de las asignaciones.
    \item \textbf{Agente viajero:} encontrar la ruta más corta que visite un 
    conjunto de ciudades exactamente una vez y regrese al punto de partida.
\end{itemize}

Suele usarse en instancias medianas que el \emph{backtracking} puro no podría 
manejar.

En general el peor caso sigue siendo exponencial, pero en la práctica las podas 
hacen que instancias razonables se resuelvan en tiempos manejables.

% ---------- Probabilísticos ----------
\section{Algoritmos probabilísticos}
\label{sec:prob}

Un algoritmo \emph{probabilístico} (también llamado aleatorizado) toma al menos 
una decisión en función de números pseudoaleatorios. Para una misma entrada, 
dos ejecuciones pueden seguir caminos distintos e incluso, en algunos casos, 
dar respuestas distintas. Suena raro al inicio, pero introducir aleatoriedad 
puede ser muy útil: a veces simplifica el algoritmo, otras lo hace 
asintóticamente más rápido y otras evita los casos peores que un algoritmo 
determinístico tendría. Es como si se introdujera entropía para escapar de la 
rigidez de los algoritmos determinísticos.

Se suelen distinguir tres familias:
\begin{itemize}
    \item \textbf{Monte Carlo.} Siempre terminan en un tiempo acotado, pero 
    pueden dar una respuesta incorrecta con cierta probabilidad. Esta 
    probabilidad se puede reducir tanto como se quiera repitiendo el algoritmo. 
    Ejemplo clásico: el test de primalidad de Miller--Rabin.
    \item \textbf{Las Vegas.} Nunca dan una respuesta incorrecta, pero su 
    tiempo de ejecución es aleatorio. Ejemplo: \emph{Quicksort} con pivote 
    aleatorio, cuyo tiempo esperado es $\Theta(n\log n)$.
    \item \textbf{Numéricos.} Producen una aproximación cuya precisión mejora 
    con el tiempo de cómputo. Ejemplo típico: la integración por Monte Carlo, 
    útil sobre todo para integrales en muchas dimensiones.
\end{itemize}

A continuación, se muestra un ejemplo de cómo se puede mejorar la probabilidad 
de acierto de un algoritmo de Monte Carlo mediante repeticiones (el cual 
pretende ser meramente ilustrativo, no se pretende el análisis riguroso):

Supóngase que se tiene un algoritmo de Monte Carlo que acierta con 
probabilidad $\geq 1/2 + \varepsilon$. Si se ejecuta $k$ veces (con $k$ impar) 
y se devuelve la respuesta más frecuente, la probabilidad de equivocarse cae 
exponencialmente con $k$.

\begin{verbatim}
    repetir_MC(x, k):
        # x = instancia, k = numero impar de repeticiones
        # MC(x) es el algoritmo Monte Carlo base (puede equivocarse)
        V = 0    # votos a "verdadero"
        F = 0    # votos a "falso"
        para i = 1, ..., k:
            si MC(x) == verdadero:
                V = V + 1
            si no:
                F = F + 1
        si V > F:
            devolver verdadero
        si no:
            devolver falso
\end{verbatim}

% ---------- Paralelos ----------
\section{Algoritmos paralelos}
\label{sec:par}

Suponga que se tiene una serie de tareas las cuales se pueden ejecutar al 
mismo tiempo, ¿convendría hacerlas secuencialmente o en paralelo?. Un algoritmo 
\emph{paralelo} toma este principio y divide el trabajo entre varios 
procesadores que actúan al mismo tiempo. Los objetivos son los de siempre que 
se piensa en paralelo: terminar antes, atacar problemas más grandes y, en 
algunos casos, ser más robustos ante fallas.

La clasificación clásica para este tipo de algoritmos (Flynn) distingue:
\begin{itemize}
    \item \textbf{SISD} (\emph{Single Instruction, Single Data}): el cómputo 
    secuencial de toda la vida.
    \item \textbf{SIMD} (\emph{Single Instruction, Multiple Data}): la misma 
    instrucción se aplica sobre datos distintos. Es lo que hacen las GPU.
    \item \textbf{MIMD} (\emph{Multiple Instruction, Multiple Data}): cada 
    procesador puede hacer cosas distintas con datos distintos. Es lo común 
    en clústeres y procesadores multinúcleo.
\end{itemize}

Dos dominan la práctica:
\begin{itemize}
    \item \textbf{Memoria compartida:} todos los procesos ven la misma memoria. 
    Se programa con hilos (\texttt{Pthreads}, \texttt{OpenMP}, hilos de Java).
    \item \textbf{Paso de mensajes:} los procesos no comparten memoria y se 
    comunican mandándose mensajes. Lo típico en clústeres y \emph{grids}; el 
    estándar es \texttt{MPI}.
\end{itemize}

Para medir el rendimiento de un algoritmo paralelo, se utilizan dos métricas 
principales. Si $T_1$ es el tiempo del algoritmo secuencial y $T_p$ el tiempo 
con $p$ procesadores, se definen:
\begin{equation}
    \text{Speedup: } S_p = \frac{T_1}{T_p}, \qquad \text{Eficiencia: } 
    E_p = \frac{S_p}{p}.
\end{equation}
La famosa \emph{ley de Amdahl} pone un techo a este \emph{speedup} en función 
de la fracción del algoritmo que no se puede paralelizar.

Por ejemplo, para sumar $n$ enteros (con $n$ potencia de 2) usando muchos 
procesadores, se organizan los datos como hojas de un árbol binario completo. 
En cada nivel los procesadores suman pares en paralelo. El resultado se obtiene 
en $\Theta(\log n)$ pasos en lugar de los $\Theta(n)$ que toma el algoritmo 
secuencial.

\begin{verbatim}
    suma_paralela(A):
        # A es un arreglo de tamano n = 2^k
        # en cada nivel se suman parejas vecinas en paralelo
        para l = 1, ..., log2(n):
            para i = 1, ..., n / 2^l   EN PARALELO:
                A[i] = A[2i - 1] + A[2i]
        devolver A[1]
\end{verbatim}

% ---------- Metaheurísticas ----------
\section{Metaheurísticas}
\label{sec:meta}

La etimología de la palabra \emph{metaheurística} proviene de las raíces 
griegas \emph{meta} (``más allá'') y \emph{heuriskein} (``descubrir''): es 
decir, una técnica que va más allá de las heurísticas tradicionales, actuando 
como una estrategia general que guía o modifica otras heurísticas para explorar 
el espacio de soluciones de manera más efectiva.

A diferencia de los algoritmos exactos, las metaheurísticas no garantizan que 
la solución encontrada sea óptima. Pero en muchos problemas reales esto no es 
un problema como tal, a cambio se consiguen muy buenas soluciones en tiempos 
razonables, que es justo lo que se necesita en la práctica.

Es especialmente útil en problemas de optimización combinatoria y continua de 
gran escala: 

\begin{itemize}
    \item \textbf{Combinatoria:} TSP, ruteo de vehículos, asignación de tareas, 
    programación entera, etc.
    \item \textbf{Continua:} optimización de funciones no convexas, ajuste de 
    parámetros en modelos complejos, etc.
    \item \textbf{Híbridos:} problemas que combinan elementos discretos y 
    continuos, como el diseño de redes o la planificación de la producción.
\end{itemize}

entre muchos otros. En general, problemas \textsf{NP}-duros donde se necesita 
una buena respuesta pronto.

Toda metaheurística tiene que balancear dos cosas:
\begin{itemize}
    \item \textbf{Exploración (diversificación):} cubrir muchas regiones 
    distintas del espacio de búsqueda para no quedarse atrapada en un óptimo 
    local.
    \item \textbf{Explotación (intensificación):} refinar la búsqueda en las 
    regiones más prometedoras para acercarse al mejor valor posible.
\end{itemize}
Lograr el balance adecuado es, en gran medida, el arte de diseñar una buena 
metaheurística.

Y se suelen dividir en dos grandes familias:
\begin{enumerate}
    \item \textbf{Basadas en una solución} (orientadas a la explotación): 
    \emph{hill climbing}, \emph{simulated annealing}, búsqueda tabú, búsqueda 
    local iterada (ILS), GRASP, búsqueda en vecindad variable (VNS).
    \item \textbf{Basadas en población} (orientadas a la exploración): 
    algoritmos genéticos y, más en general, algoritmos evolutivos, estrategias 
    de evolución, optimización por enjambre de partículas (PSO), colonias de 
    hormigas (ACO), búsqueda dispersa, sistemas inmunes artificiales.
\end{enumerate}

Algo interesante es que el análisis asintótico clásico no dice mucho 
de las metaheurísticas, y esto es esencial en la práctica porque obtienen una
solución muy buena en tiempos razonables. Por eso suelen evaluarse 
experimentalmente, sobre instancias de prueba, comparando estadísticamente 
con otros métodos.

% --------------------------- COMPARATIVO -----------------------------
\section{Análisis comparativo}
\label{sec:comparativo}

A continuación se muestra la tabla~\ref{tab:comparativa} que resume las 
propiedades principales de cada técnica. Se contrastan complejidad temporal 
típica, uso de memoria, si garantizan o no la solución óptima y dónde suelen 
usarse.

\begin{table}[ht]
\centering
\scriptsize
\renewcommand{\arraystretch}{1.2}
\caption{Comparativa de las principales técnicas de diseño algorítmico.}
\label{tab:comparativa}
\begin{tabular}{|p{2.2cm}|p{3.2cm}|p{2.3cm}|p{2.2cm}|p{4cm}|}
\hline
\textbf{Técnica} & \textbf{Complejidad} & \textbf{Memoria} & \textbf{¿Óptimo?} & \textbf{Aplicaciones} \\
\hline
Divide y vencerás       & $\mathcal{O}(n\log n)$ a $\mathcal{O}(n^c)$ & $\mathcal{O}(\log n)$ (pila)   & Sí                & Mergesort, búsqueda binaria, Strassen \\
\hline
Voraces                 & $\mathcal{O}(n\log n)$ típico               & Baja                            & Solo con subestr. óptima & MST, Huffman, planificación \\
\hline
Programación dinámica   & Polinómica o pseudopolinómica               & $\Theta(\text{tabla})$, alto    & Sí                & Mochila 0-1, distancia de edición \\
\hline
Backtracking            & Exponencial (peor caso)                     & $\mathcal{O}(n)$ (pila)        & Sí (exhaustivo)   & $n$-reinas, coloración \\
\hline
Ramif. y poda           & Exponencial; mejor con poda                 & Variable                        & Sí                & Mochila 0-1, TSP \\
\hline
Probabilísticos         & Las Vegas: variable; Monte Carlo: acotada   & Baja                            & Probabilística    & Miller--Rabin, Quicksort aleatorio \\
\hline
Paralelos               & Reducción por factor $\leq p$ (Amdahl)      & Depende del modelo              & Igual al base     & Cómputo científico, IA \\
\hline
Metaheurísticas         & Según presupuesto                           & Variable                        & No                & Optimización \textsf{NP}-difícil \\
\hline
\end{tabular}
\end{table}

\subsection{¿Cómo elegir?}
Después de haber visto todos los algoritmos y una breve comparación sobre 
ellos, surge la pregunta natural, ¿cómo elegir la técnica adecuada? 
Bueno, la elección de la técnica depende, en la práctica, de cuatro cosas:

\begin{enumerate}
    \item \textbf{Naturaleza del problema.} Si admite subestructura óptima, 
    voraces o programación dinámica son buenos candidatos. Si requiere explorar 
    combinaciones, \emph{backtracking} o ramificación y acotación.
    \item \textbf{Tamaño de la entrada.} Para entradas pequeñas, los algoritmos 
    exactos exponenciales son perfectamente usables. Para entradas grandes en 
    problemas difíciles, prácticamente no queda más opción que las 
    metaheurísticas.
    \item \textbf{Recursos disponibles.} La programación dinámica gasta mucha 
    memoria, el paralelismo requiere \emph{hardware} adecuado, las 
    metaheurísticas requieren tiempo de ajuste de parámetros.
    \item \textbf{Garantías requeridas.} Si se necesita el óptimo exacto, los 
    métodos aproximados quedan descartados. Si basta con ``muy bueno'', se abre 
    todo el abanico.
\end{enumerate}

Aunque es solo teoría, en la práctica casi nunca se usa una técnica pura. Lo 
común es combinar: por ejemplo, una metaheurística que arranca con una 
heurística voraz para tener una buena solución inicial, usa programación 
dinámica para resolver exactamente algunos subproblemas y aprovecha varios 
núcleos en paralelo. La hibridación es más la regla que la excepción.

% --------------------------- CONCLUSIONES ----------------------------
\section{Conclusiones}
\label{sec:conclusiones}

A lo largo del trabajo se revisaron las principales familias de algoritmos. 
Las ideas más importantes que conviene llevarse son:

Primero, los algoritmos se pueden clasificar por más de un criterio: por su 
comportamiento (determinístico, no determinístico, probabilístico), por la 
estructura del cómputo (secuencial o paralelo) o por la técnica de diseño 
(divide y vencerás, voraz, programación dinámica, vuelta atrás, ramificación 
y acotación). Estas dimensiones no se excluyen: un mismo algoritmo puede 
pertenecer a varias categorías a la vez.

Y segundo, no hay una técnica universalmente mejor. La elección depende del 
problema concreto, del tamaño de las entradas, de los recursos disponibles y 
de qué garantías se necesitan sobre la solución.

% --------------------------- BIBLIOGRAFÍA ----------------------------
\section{Bibliografía}
\label{sec:bibliografia}

\begin{enumerate}
    \item Cormen, T. H., Leiserson, C. E., Rivest, R. L. y Stein, C. \emph{Introduction to Algorithms}. 3.ª ed. The MIT Press, Cambridge, Massachusetts, 2009.
    \item Brassard, G. y Bratley, P. \emph{Fundamentos de Algoritmia}. Prentice Hall, Madrid, 1997. Traducción del original publicado por el Département d'informatique et de recherche opérationnelle, Université de Montréal.
    \item Sipser, M. \emph{Introduction to the Theory of Computation}. 3.\textsuperscript{rd} ed. Cengage Learning, Boston, Massachusetts, 2012.
    \item Talbi, E.-G. \emph{Metaheuristics: From Design to Implementation}. John Wiley \& Sons, Hoboken, Nueva Jersey, 2009. University of Lille - CNRS - INRIA.
\end{enumerate}

\end{document}